\section{Finite frame rings and Cartier transfers}
\label{sec:rings}

To compare coefficient cohomology with compact supports, we need transfers
whose transpose is the actual Frobenius on functions of the finite covers. This
section constructs the finite frame rings and lift torsors, then normalizes
their Cartier transfers to obtain that map in Lemma~\ref{lem:cartier-transfer}.

Fix integers $1\leq t<n<p-1$, and put $m=n-t$.  Choose a finite
field $k$ containing the fields of definition of the Honda endomorphisms
at the heights at most $n$.  Its degree over $\Fp$ may be chosen prime
to $p$: a multiple of $\operatorname{lcm}(1,\ldots,n)$ suffices for
these endomorphism fields.  Further finite extensions used below are
chosen with the same property.  Write
\[
 A=k[[x_t,\ldots,x_{n-1}]],\qquad x=x_t,\qquad B=A[1/x].
\]
Here $A$ is the height-$t$ quotient of the characteristic-$p$
Lubin--Tate deformation ring of the height-$n$ Honda group.  Write
$\mathcal G_A$ for its formal group and
$P_n=\mathcal O_{D_n}^{\times}$ for the ordinary height-$n$
stabilizer.  The action on a function induced by a left action on
points is always inverse substitution.

We choose a coordinate on $\mathcal G_A$ lifting the fixed Honda
coordinate and having the height congruences
\begin{equation}\label{rings:height-congruences}
 [p](T)\equiv x_iT^{p^i}
 \pmod{(x_t,\ldots,x_{i-1},T^{p^i+1})},
 \qquad t\leq i\leq n,
 \quad x_n=1.
\end{equation}
The same choice gives the full special-fiber identity
\[
 [p](T)\equiv T^{p^n}\pmod{(x_t,\ldots,x_{n-1})}
\]
as an identity of formal series.  Moreover $[p](T)$ factors through
$T^{p^t}$ over $A$.
The group used on $B$ is the algebraic $p$-divisible group obtained
from the finite kernels over $A$.  More explicitly, Weierstrass
preparation makes
\[
 A[[T]]/([p^l](T))
\]
a finite free $A$-algebra of rank $p^{nl}$.  Its torsion coordinate is
topologically nilpotent before localization, so the formal addition
and multiplication series define algebraic group operations on this
finite algebra.  Preparation for $[p](T)-Z$ gives the finite flat
transition maps.  We localize these finite group schemes at $x$.
Thus the connected--\'etale sequence on $B$ retains the entire finite
kernel, including its infinitesimal connected part.

\subsection{Labeled rings before localization}

For $l\geq1$, consider $m$ points of $\mathcal G_A[p^l]$ whose
images form a basis of the \'etale quotient at the generic point of
$A$.  Let $\widetilde C_l$ be the reduced closure of this generic
lifted-basis locus in the finite scheme of $m$ torsion points, and
write $z_{1,l},\ldots,z_{m,l}$ for the lifted coordinates.  Set
\[
 w_{i,l}=z_{i,l}^{p^{tl}},\qquad
 C_l=A[w_{1,l},\ldots,w_{m,l}]\subseteq\widetilde C_l.
\]
These finite $A$-algebras are complete.  A reduced closure is used
only to establish their coordinates; the proposition also identifies
the unreduced scheme that will be used for descent.
For the boundary convention $t=n$, the empty labeled ring is $k$.

\begin{proposition}[Finite labeled rings]\label{prop:finite-rings}
There are compatible identifications
\begin{equation}\label{rings:regular-coordinates}
 C_l=k[[w_{1,l},\ldots,w_{m,l}]],\qquad
 \widetilde C_l=k[[z_{1,l},\ldots,z_{m,l}]],\qquad
 w_{i,l}=z_{i,l}^{p^{tl}}.
\end{equation}
Both rings are finite free over $A$.  On the exact-height locus,
$C_l[1/x]$ is the finite \'etale algebra of bases of the \'etale
$p^l$-torsion.  Over this algebra the scheme of all lifts of its
universal basis to $\mathcal G_A[p^l]$ has algebra
\[
 \widetilde C_l[1/x]=(C_l[1/x])^{1/p^{tl}},
\]
finite locally free of rank $p^{tlm}$.

The transition maps are finite and injective, and for $j\leq l$
\begin{equation}\label{rings:transition-powers}
 z_{i,j}\in
 k[[z_{1,l}^{p^{t(l-j)}},\ldots,z_{m,l}^{p^{t(l-j)}}]].
\end{equation}
\end{proposition}

\begin{proof}
The algebra $\widetilde C_l$ is finite and reduced, and each of its
irreducible components dominates $\Spec A$ by its definition as a
generic closure.  It has dimension $m$.  Above the closed point of
$A$ all torsion coordinates are nilpotent, so there is a unique
closed point and its residue field is $k$.  Its maximal ideal is
generated by the height parameters and the $z_{i,l}$.

Put $a_i=[p^{l-1}](z_{i,l})$, and for $v\in\Fp^m$ let $a_v$ be
the corresponding formal-group linear combination of the $a_i$.
If $W_p(T)$ is the distinguished polynomial for $[p](T)$, then
\begin{equation}\label{rings:basis-polynomial}
 W_p(T)=\prod_{v\in\Fp^m}(T-a_v)^{p^t}
       =\prod_{v\in\Fp^m}(T^{p^t}-a_v^{p^t}).
\end{equation}
Indeed this is the factorization over every geometric generic
field: the distinct \'etale points are indexed by the labels and
each has connected multiplicity $p^t$.  Both sides are monic of
degree $p^n$, so generic density in the reduced algebra proves the
identity there.

There is a series $h_l$ with
$[p^{l-1}](T)=h_l(T^{p^{t(l-1)}})$.  Consequently
$a_i^{p^t}$ is a power series in $w_{i,l}$ with coefficients in
$A$ and zero constant term.  Taking formal-group linear
combinations gives the same assertion for every $a_v^{p^t}$.
Thus \eqref{rings:basis-polynomial} is an identity over $C_l$,
and modulo $(w_{1,l},\ldots,w_{m,l})$ it becomes $T^{p^n}$.
The preparation unit for $[p](T)$ is defined over $A$.  All
coefficients of $[p](T)$ of degree below $p^n$ therefore vanish in this
quotient, and \eqref{rings:height-congruences} successively gives
\[
 x_t=x_{t+1}=\cdots=x_{n-1}=0.
\]
The maximal ideal of $C_l$ is generated by its $m$ displayed
$w$-coordinates.  Since $C_l$ is integral over $A$, its dimension
is $m$.  It is a regular local ring, and the surjection from
$k[[W_1,\ldots,W_m]]$ taking $W_i$ to $w_{i,l}$ is an
isomorphism: a nonzero kernel would lower the dimension of this
power-series domain.  The identical maximal-ideal argument with
the $z_{i,l}$ gives the second identification in
\eqref{rings:regular-coordinates}, including the stated
inclusion.  These finite regular rings are Cohen--Macaulay of
dimension $m$; the parameters of $A$ form systems of parameters
on them.  Their Koszul complexes are exact in positive degree,
so the finite modules are flat, hence free, over the local ring
$A$.

We next check the finite schemes after localization.  Write
\[
 [p^l](T)=g_l(T^{p^{tl}}),\qquad
 g_l(S)=U_l(S)V_l(S),
\]
where $U_l$ is a unit and $V_l$ is distinguished of degree
$p^{ml}$.  The linear coefficient of $g_l$ is a unit times
$x^{1+p^t+\cdots+p^{t(l-1)}}$.  The invariant-differential
identity for this Frobenius-divided homomorphism expresses
$g_l'(S)$ as that coefficient times a quotient of unit series.
Evaluating in the prepared finite algebra shows that $V_l'$ is
a unit after $x$ is inverted.  Thus the Frobenius-divided finite
group is \'etale there, of rank $p^{ml}$.  Its kernel quotient
has removed exactly the connected kernel of rank $p^{tl}$.

The scheme of bases of this \'etale group is finite \'etale over
$B$.  The universal labeled points give a surjection from its
algebra to $C_l[1/x]$.  It is an isomorphism generically, since
the generic closure included every basis.  Its kernel is zero:
the source is flat over the domain $B$, whereas a kernel
vanishing generically would be $B$-torsion.  This identifies the
basis algebra.

Over this basis algebra, the scheme of lifts of the $m$ basis
vectors is a product of torsors for the connected finite kernel.
It is finite locally free of rank $p^{tlm}$.  Its algebra
surjects onto $\widetilde C_l[1/x]$, which is generated by the
same lifted coordinates.  The latter is free over $C_l[1/x]$
with basis
\[
 \prod_{i=1}^m z_{i,l}^{e_i},\qquad 0\leq e_i<p^{tl}.
\]
The two ranks are equal, so the surjection is an isomorphism.
This is the assertion about the full lift scheme, rather than
its reduction.

A basis at level $j$ lifts to a basis at level $l$ over an
extension of every geometric generic field.  Hence the
transition maps on the reduced closures are injective.  They
are finite because both rings are finite over $A$.  Since
$z_{i,j}=[p^{l-j}](z_{i,l})$, factorization through Frobenius
and \eqref{rings:regular-coordinates} give
\eqref{rings:transition-powers}.

\end{proof}

The underlying regular-quotient and \'etale-coordinate
construction is also described at arbitrary height in
\cite[\S\S3--4, Proposition~4.1]{Strauch}.  The explicit
rank comparison above records the additional finite-flat
statement needed for our translation torsors.

The boundary calculation also needs a presentation that survives
nilpotent base change.  We establish that presentation before
identifying the locus on which a label vanishes.

\begin{lemma}[An integral full-section presentation]
\label{rings:full-sections-presentation}
Let $H$ be the quotient of $\mathcal G_A[p]$ by its relative
$p^t$-Frobenius kernel.  The ring $C_1$ represents full labeled
section tuples $\Fp^m\to H$, including over nilpotent test
rings.  Here ``full'' means equality of the finite Cartier
divisors, or equivalently, after every base change,
equality of the norm of every function with its product
over the labeled sections.
\end{lemma}

\begin{proof}
The quotient $H$ has the monogenic presentation obtained by
dividing the prepared $p$-series by relative Frobenius.  The
full-section condition is a finite set of coefficient
identities for its monic polynomial.  Indeed equality for the
coordinate implies equality of the characteristic polynomial
for every polynomial in that coordinate by substitution, and
every function in the finite algebra has such a representative.
This description commutes with arbitrary base change.

Let $Q$ be the resulting finite complete $A$-algebra of labeled
group-homomorphism tuples.  It has one closed point, because
at the closed point of $A$ the connected finite group has only
the zero section over a field.  If all $m$ labeled coordinates
are set equal to zero, the full-section identity makes the
divided distinguished polynomial $T^{p^m}$.  Applying
\eqref{rings:height-congruences} before this division
successively kills all the height parameters.  Thus the
maximal ideal of $Q$ has at most $m$ generators.  There is a
surjection $Q\twoheadrightarrow C_1$ induced by the generic
full tuple.  Since $C_1$ has dimension $m$, the dimension and
embedding dimension of $Q$ are both $m$.  Hence $Q$ is regular.
A quotient of this regular local domain of the same dimension
has zero kernel, proving $Q=C_1$.

All equations in this argument are equations in finite
algebras formed before localization.  Their subsequent use
on an \'etale component therefore involves polynomial
operations in a finite algebra; it does not require substituting
an order-one element into an unevaluated formal series.
\end{proof}

\begin{corollary}[Vanishing labels and higher starting height]
\label{rings:boundary-labels}
At level one, for each nonzero label $v\in\Fp^m$ there is a
coordinate $w_v$ such that
\begin{equation}\label{rings:height-divisor}
 x=\text{a unit}\cdot\prod_{0\ne v\in\Fp^m}w_v.
\end{equation}
Each ideal $(w_v)$ is $P_n$-stable.  Its conormal is a power of the
invariant coordinate line.  For a primitive label, the quotient by
$(w_v)$ is the starting-height-$(t+1)$ level-one ring with $p$-th
roots of its parameters adjoined.  The same assertion can be
iterated for any independent set of dead labels.
\end{corollary}

\begin{proof}
At level one let $a_v$ be the formal-group linear combination of
$z_{1,1},\ldots,z_{m,1}$ with label $v\in\Fp^m$, and put
$w_v=a_v^{p^t}$.  The coefficient of $T^{p^t}$
in \eqref{rings:basis-polynomial}, together with the
preparation unit, gives \eqref{rings:height-divisor}.
Transport by $P_n$ is a change of formal coordinate with
invertible linear coefficient.  Its action on $w_v$ is $w_v$
times a unit, proving stability of $(w_v)$.  On the conormal
only that linear coefficient remains, raised to the $p^t$-th
power.  Equivalently, this is the invariant cotangent line of
the Frobenius-divided group restricted to its zero section.
This also describes the action without a choice of generator
for the line.

Change the labels so that $v$ is the first basis vector, and
put $R=C_1/(w_1)=k[[w_2,\ldots,w_m]]$.
Write $[p](T)=g_t(T^{p^t})$ and let $Q_t$ be the distinguished
polynomial of $g_t$.  Formula~\eqref{rings:height-divisor}
gives $x_t=0$ in $R$.  The labels in each coset of $\Fp v$
then give the same point, so
\[
 Q_t(Y)=\prod_{\bar v\in\Fp^{m-1}}
              (Y^p-w_{\bar v}^{p}).
\]
Over $A/(x_t)$ we have $g_t(Y)=g_{t+1}(Y^p)$, and uniqueness
of preparation gives $Q_t(Y)=Q_{t+1}(Y^p)$.  Therefore
\[
 Q_{t+1}(Z)=\prod_{\bar v\in\Fp^{m-1}}(Z-w_{\bar v}^p).
\]
Coefficient Frobenius identifies these labels with the
formal-group combinations of $w_2^p,\ldots,w_m^p$ in the
next Frobenius twist.  By Lemma~\ref{rings:full-sections-presentation},
the displayed full-section identity gives a map from the
starting-height-$(t+1)$ labeled ring, whose labeled
coordinates map to $w_2^p,\ldots,w_m^p$.
In the power-series coordinates of
Proposition~\ref{prop:finite-rings}, this map is injective with image
$k[[w_2^p,\ldots,w_m^p]]=R^p$, since $k$ is perfect.
Thus $R$ is exactly its $p$-th-root extension, with the
indicated $A/(x_t)$-structure.  Repeating the calculation
for $h$ independent dead labels sends the higher-height
labeled coordinates to $w_{h+i}^{p^h}$ and gives the corresponding
$p^h$-th-root extension.
\end{proof}

\subsection{Connected markings and compatible lifts}

Let $\Gamma_t$ be the height-$t$ Honda group over $k$.
A connected marking is an isomorphism
\[
 \alpha:\Gamma_t\longrightarrow \mathcal G_B^0
\]
of formal groups, and an \'etale marking is an integral basis
$i:\Zp^m\longrightarrow T_p\mathcal G_B^{\mathrm{et}}$.
Their joint algebraic tower has structure group
\[
 U_0=P_t\times\GL_m(\Zp),\qquad
 P_t=\mathcal O_{D_t}^{\times}.
\]
We use the left action
$(h,a)(\alpha,i)=(\alpha h^{-1},ia^{-1})$.
For an open subgroup $U\subseteq U_0$, write $B_U$ for the
corresponding finite frame algebra.  Write
$\mathcal B=\colim_U B_U$ for the entire joint algebraic
frame ring.  No completion of this union is part of this
notation.

\begin{proposition}[Frame and lift torsors]\label{prop:frame-torsor}
The rings $B_U$ form a tower of finite \'etale $B$-algebras.
After passing to normal open levels, the maps in the tower
are finite \'etale torsors for the indicated finite deck
groups.  The tower is $P_n$-equivariant, and $P_n$ commutes
with its $U_0$-action.

Over $\mathcal B$, the scheme of lifts of the compatible
\'etale basis is an affine faithfully flat torsor for
\begin{equation}\label{rings:translation-group}
 T=(T_p\Gamma_t)^m,
 \qquad
 T_p\Gamma_t=\varprojlim_{[p]}\Gamma_t[p^r].
\end{equation}
For every $r\geq1$ its finite quotient has group
$T/[p]^rT=\Gamma_t[p^r]^m$ and algebra
\begin{equation}\label{rings:root-torsor}
 \mathcal B^{1/p^{tr}}.
\end{equation}
Thus the full lift algebra is the algebraic root union
$\mathcal B^{\mathrm{perf}}$.
The group $P_n$ acts trivially on the translation parameters.
The action of $(h,a)\in U_0$ on a translation tuple $\tau$ is
\begin{equation}\label{rings:auxiliary-action}
 \tau\longmapsto h\tau a^{-1}.
\end{equation}
All finite torsor constructions and all their morphisms
descend to finite frame levels once their finite marking
data are included.
\end{proposition}

\begin{proof}
Over $B$ the lower height coefficients vanish and the
height-$t$ coefficient is invertible.  The full isomorphism
ring between this law and $\Gamma_t$ is therefore a filtered
union of finite \'etale extensions, by the exact-height
isomorphism theorem \cite[Lecture~14, Theorem~1]{LurieFormalGroups}.
For clarity, its finite stages record coefficients of a full
isomorphism that extend to the next required equations.
The first coefficient $c$ satisfies
$c^{p^t-1}=x^{-1}$ up to the chosen leading unit.  The
subsequent free coefficients satisfy monic additive equations
with invertible linear coefficient; the other coefficients
are uniquely determined.  The equation constraining a
coefficient of degree $p^j$ can occur at degree $p^{t+j}$.
Thus these are not the schemes of arbitrary finite-bud
isomorphisms.  Any two full markings differ uniquely by
$P_t$, and quotienting by an open subgroup gives its finite
\'etale torsor stages.

The \'etale Tate basis tower is finite \'etale at each
level, as proved in Proposition~\ref{prop:finite-rings}.
Taking the product of the two frame torsors proves the
first assertion.  Universal deformation transport carries
both markings along the same isomorphism; it commutes with
changing the markings.  This proves the asserted $P_n$
equivariance.

At finite level $r$, two lifts of the marked \'etale basis
differ by exactly one element of $\Gamma_t[p^r]^m$.
Proposition~\ref{prop:finite-rings} identifies the lift
algebra over the \'etale basis cover with its $p^{tr}$-th
root ring.  Relative Frobenius commutes with finite \'etale
base change, so the same identification holds after any
connected marking base change.  It is canonical: each
element of the lift algebra is the unique $p^{tr}$-th root
of its power in the base.  The identity may be checked on
the faithfully flat basis cover and then descended.

The finite torsor maps are faithfully flat and the
transition on lifts is multiplication by $p$.  Taking
their affine inverse limit gives the torsor identity and
faithful flatness for \eqref{rings:translation-group}.
The kernel of the projection $T\to\Gamma_t[p^r]^m$ is
$[p]^rT$: shifting a compatible sequence by $r$ constructs
its unique preimage under $[p]^r$.  This proves both the
subgroup assertion and \eqref{rings:root-torsor}.

Finally, transport by $P_n$ carries the connected marking,
the \'etale basis, and its lift simultaneously, so it fixes
their difference parameter in the constant Honda group.
Changing the markings by $(h,a)$ instead changes that
parameter by \eqref{rings:auxiliary-action}.  Each finite
kernel and each finite diagram uses finitely many
coefficients of the connected marking and a finite
\'etale level, proving finite-stage descent.
\end{proof}

\Needspace{8\baselineskip}
\begin{definition}[Coefficient and cochain convention]
\label{rings:bounded-cochains}
For a finite projective module over a finite frame algebra
we use continuous cochains with bounded poles.  On a compact
profinite source this means that one power of $x$ places the
entire image in a finite power-series lattice, where the
cochain is continuous for the parameter-adic topology.
For a union of frame, root, or associated-bundle coefficients,
one common finite stage is required on each compact source.
Equivalently the cochain complex of such a union is the
filtered colimit of its finite-stage bounded-pole cochain
complexes.  These conventions also apply with an additional
compact profinite parameter.
\end{definition}

Finite free coordinate formulas, module splittings, and
finite \'etale descent preserve this convention.  A fixed
finite diagram can have a fixed pole loss, but its formulas
require only finitely many coefficients and denominators.
The compact-lattice estimates and continuous cochain
exactness used later will be established with these
coefficient topologies.

\subsection{The dual distribution ring}

The distribution ring organizes the translation representations at each
finite root level. Its identifications across Frobenius levels will be used in
Lemma~\ref{lem:cartier-transfer} to normalize the transfer and identify its
transpose with the actual Frobenius on cover functions.

Suppose first that $t>1$, and put $d=m(t-1)$.
The category of rational representations of the affine
group scheme $T$ is the category in which every vector
belongs to a representation of a finite quotient of $T$.
It is not the representation category of its geometric
points.  Finite Cartier duality identifies its completed
distribution ring with
\begin{equation}\label{rings:distribution-ring}
 \Lambda=\varprojlim_r\mathcal O(\Gamma_t[p^r]^m)^\vee
     =\mathcal O((\widehat{\Gamma_t^\vee})^m)
     \simeq k[[D_1,\ldots,D_d]].
\end{equation}
Here $(-)^\vee$ on a finite vector space is its $k$-linear
dual.  Rational representations correspond to discrete
modules locally killed by powers of the augmentation ideal.
The displayed parameters need not linearize the auxiliary
action.

\begin{lemma}[Divisible Frobenius levels]\label{rings:dual-distribution}
There is a positive integer $r_0$, divisible by $t-1$, such
that for $r=ar_0$, $a\geq1$, the two numbers
\begin{equation}\label{rings:two-frobenius-indices}
 Q=p^{tr},\qquad q=p^{tr/(t-1)}
\end{equation}
are Frobenius powers fixing $k$.  At this level
\[
 \mathcal O((\Gamma_t^\vee[p^r])^m)
   =\Lambda_q:=\Lambda/(D_1^q,\ldots,D_d^q),
 \qquad
 \dim_k\Lambda_q=q^d=Q^m.
\]
The embedding $[p]^r:T\hookrightarrow T$ corresponds on
distributions to
\[
 \psi_q:\Lambda\longrightarrow\Lambda,
 \qquad D_j\longmapsto D_j^q,\quad c\longmapsto c\ (c\in k).
\]
If $N_q=\mathcal O(T/[p]^rT)$ is the regular function
representation, finite Hopf duality gives a normalized
equivariant identification $N_q\simeq\Lambda_q$ as
translation representations.  The auxiliary action on
each finite quotient factors through a finite group.
\end{lemma}

\begin{proof}
The dual of the dimension-one, height-$t$ Honda group has
dimension $t-1$ and height $t$.  Since $t>1$, it is connected;
its formal group is smooth of dimension $t-1$.  This proves
\eqref{rings:distribution-ring}.  The Honda relation is
$F^t=[p]$ on $\Gamma_t$.  Duality gives $V^t=[p]$ on
$\Gamma_t^\vee$, and $FV=[p]$ then gives
\[
 F^t=[p]^{t-1}\quad\hbox{on }\Gamma_t^\vee.
\]
Cancellation here is cancellation of an isogeny of the
divisible group.  Thus, when $r$ is divisible by $t-1$,
$[p]^r$ on the dual is Frobenius of order $tr/(t-1)$.
Choose $r_0$ also so that this Frobenius and Frobenius of
order $tr_0$ fix $k$.  Pullback on any system of formal
parameters is consequently the stated $q$-power map.
Its kernel has the displayed truncated power-series ring.
The length identity follows directly from $d=m(t-1)$.

For a finite commutative group scheme $G$, its function
algebra and the dual Hopf algebra are related by the
perfect Frobenius pairing
\[
 (f,g)\longmapsto\lambda(fg),
\]
where $\lambda$ is a nonzero invariant integral.  In our
Frobenius kernels, a coordinate description of $\lambda$
is the highest-monomial coefficient after multiplication
by the normalized invariant volume density.  The pairing
is perfect: its associated graded pairing in augmentation
coordinates pairs complementary monomials.  More explicitly,
let $v$ be the highest monomial in the original group
coordinates.  The map $\Lambda_q\to N_q$, $a\mapsto a\cdot v$,
is an isomorphism of translation representations.  Indeed
the socle integral of the local Gorenstein algebra
$\Lambda_q$ acts nontrivially on $v$.  A nonzero kernel
ideal would meet that socle, a contradiction, and the two
dimensions agree.  Under an auxiliary change of group
coordinates, $v$ changes by the tangent determinant
raised to the $(Q-1)$-st power; higher coordinate terms
have too large total degree to contribute to the socle.
This character is trivial, because its values lie in
$k^\times$ and $Q$-power fixes $k$.  The displayed
regular-module identification is therefore equivariant.
With its socle normalization $\lambda(v)=1$, it satisfies
\[
 \lambda(a\cdot v)=\varepsilon(a)\lambda(v)=\varepsilon(a),
\]
where $\varepsilon$ is the distribution augmentation.
Thus normalized integration is exactly the quotient
$\Lambda_q\to k$ under this identification.
The
finite-dimensional algebra is a finite set, since $k$
is finite.  Continuity of the marking action implies
that its action factors through a finite quotient.
\end{proof}

When $t=1$ we instead take $\Gamma_1=\widehat{\mathbb G}_m$.
Its Cartier dual is \'etale.  The translation group $T$ is
diagonalizable with character group
$(\Qp/\Zp)^m$.  A rational representation is the direct
sum of its character spaces, and invariants are the
zero-character summand.  Invariants are consequently
exact.  The transfer at a finite level is the
constant-character projection, normalized to fix constants.
We set $d=0$ in this case and use no power-series
distribution ring or higher translation Ext.  The
function-root exponent remains $Q=p^r$.

\subsection{Powering and natural finite diagrams}

Write $\mathcal E=\mathcal B^{\mathrm{perf}}$ for the
algebraic lift union.  For a finite translation
representation $M$, its associated coefficient module is
\[
 \mathcal V(M)=(\mathcal E\otimes_k M)^T.
\]
It is obtained by finite flat descent from a finite lift
torsor as soon as $M$ factors through a finite quotient.
In particular it is finite projective over $\mathcal B$
and descends to a finite frame stage.  A compatible
auxiliary action on $M$ is retained in this construction.

\begin{lemma}[Powering on associated diagrams]\label{rings:coinduction}
Let $r$ be one of the chosen levels and put $H=[p]^rT$.
Transport a finite $T$-representation $M$ to $H$ using
$[p]^r:T\simeq H$, and then coinduce it to $T$.  This is
an exact functor, denoted $\mathcal C_r$.
There is a natural $P_n$- and auxiliary-equivariant
identification of the associated coefficient diagram for
$\mathcal C_rM$ with the $Q$-root version of the diagram
for $M$.  Powering by $Q$ identifies the two diagrams
$k$-linearly and Frobenius-semilinearly over the base:
\begin{equation}\label{rings:power-scalar-rule}
 \Phi_Q(cf)=c^Q\Phi_Q(f),\qquad \Phi_Q(f)=f^Q.
\end{equation}
The assertion holds for every morphism and every finite
diagram, and after invariants in any compatible open
frame subgroup.  It does not require trivializing the
auxiliary action on all coinduced modules simultaneously.
\end{lemma}

\begin{proof}
Let $\rho$ be the coaction on $\mathcal E$ and $\rho_H$
its restriction to $H$, transported to $T$.  On
$\mathcal O(T)$, substitution along $[p]^r$ is the
$Q$-power map, because $[p]^r=F^{tr}$ on the original
Honda group.  Taking powers in the coaction identity
gives
\begin{equation}\label{rings:coaction-powering}
 \rho\Phi_Q=(\Phi_Q\otimes1)\rho_H.
\end{equation}
The algebra $\mathcal E$ is perfect, so $\Phi_Q$ is
bijective.  It is $k$-linear by the choice of $r$, and
has the scalar behavior \eqref{rings:power-scalar-rule}.
It identifies $\mathcal E^H=\mathcal B^{1/Q}$ with
$\mathcal E^T=\mathcal B$.

Finite induction and coinduction for $H\subseteq T$
are exact.  For $t>1$ this can be seen directly from
\eqref{rings:distribution-ring}: $\Lambda$ is finite
free over $\psi_q(\Lambda)$, with the monomials of
exponents less than $q$ as a basis, so both tensor
induction and Hom coinduction are exact.  After
forgetting auxiliary actions their usual Frobenius
pairing also identifies these two functors.  The
equivariance needed here will follow directly from
the coaction identity and adjunction below.
For $t=1$ the same statement
is the corresponding finite character decomposition.
The descent identity for coinduction is
\[
 \mathcal V(\mathcal C_rM)
   \simeq(\mathcal E\otimes_k M)^H.
\]
Tensoring \eqref{rings:coaction-powering} with $M$ and
taking invariants gives the claimed comparison.

Every marking automorphism commutes with $[p]^r$,
and every characteristic-$p$ ring automorphism
commutes with $Q$-powering.  Hence these identities
are equivariant before taking invariants.  They are
also natural before invariants, so they apply to
entire finite diagrams, not just to their dimensions
or individual terms.  Finite torsor formulas use
finitely many marking coefficients and denominators.
They therefore descend, with their natural actions,
and respect Definition~\ref{rings:bounded-cochains}.
\end{proof}

\subsection{Invariant volumes and the transfer transpose}

The dimension of the finite lift group
$\Gamma_t[p^r]^m$ in tangent directions is $m$, whereas
the dimension of its formal Cartier dual is $d$.
The volume in this subsection has degree $m$.
Choose normalized invariant coordinate volumes on the
Honda factors.  On a sufficiently divisible finite
lift torsor, their product gives a generator of the
relative top differentials.  In the root coordinates
of Proposition~\ref{prop:finite-rings}, the base
consists of $Q$-th powers.  Relative and absolute
continuous differentials therefore agree.  Transport
this generator by the abstract ring isomorphism
$f\mapsto f^Q$ from the root ring to the original
ring, and call the resulting top form $\eta$.
This transport is transport through an isomorphism
of rings with renamed root coordinates; it is not
the zero pullback of differentials under an absolute
Frobenius morphism.
It gives compatible frames on subsequent levels.  To
check this, trivialize the relevant finite torsor
diagram faithfully flatly.  The subgroup map is
$Q$-power on the original Honda coordinates, and
the normalized invariant-volume coefficients are
fixed by this power on $k$.  Power-transport is
therefore the normalized group volume at the
preceding stage.  Equality descends along the
trivialization.  Thus one choice of $\eta$ is
used throughout the stationary tower.

The construction uses only finite torsor and marking
data.  Thus $\eta$ is a generator at a sufficiently
deep finite frame level.  On the full joint tower it
is $P_n$-invariant, since $P_n$ fixes the translation
parameters.  Its auxiliary line is
\begin{equation}\label{rings:volume-character}
 \langle\eta\rangle
    =\det\Lie(\Gamma_t^m)^*.
\end{equation}
In particular its $\GL_m$ weight is $+1$ in each
diagonal row, by \eqref{rings:auxiliary-action}.
The character is finite over $k$; it becomes trivial
on a sufficiently small open frame subgroup.

\Needspace{9\baselineskip}
\begin{lemma}[Cartier transfer and its transpose]
\label{lem:cartier-transfer}
Fix one sufficiently divisible step with function-root
exponent $Q$, and use the volume just constructed.
The quotient transfers of the lift torsors can be
normalized compatibly so that, after powering to a
fixed finite frame ring, their first step is
\begin{equation}\label{rings:stationary-c}
 S(f)=\frac{\mathscr C_Q(f\eta)}{\eta},
\end{equation}
where $\mathscr C_Q$ is Cartier trace on continuous
top differentials.  The subsequent stationary
transfers are $S^a$.  In particular
$S(c^Qf)=cS(f)$.

These transfers are finite Hopf integration on the
torsor, not ordinary field trace.  They commute
with $P_n$, with their natural auxiliary actions,
and with finite \'etale change of frame level.
Their transpose under the residue pairing is
positive Frobenius on dual coefficient functions,
with the one volume line \eqref{rings:volume-character}
retained.  Before stationary identification it is
the inclusion into the next root coefficient ring.
For a finite \'etale cover or its associated
coefficient module, this is the actual Frobenius
on the cover functions and their trace duals.
All statements hold for finite diagrams in
Lemma~\ref{rings:coinduction}.
\end{lemma}

\begin{proof}
We first verify both the normalization and the absence
of a base-dependent scalar.  On a root chart write
\[
 R_0=k[[w_1,\ldots,w_m]],\qquad
 R_1=k[[z_1,\ldots,z_m]],\qquad w_i=z_i^Q,
\]
with the relevant height parameter inverted and with
finite \'etale scalar extensions permitted.  Write
the invariant torsor volume as
\[
 \eta_1=a(z)\,dz_1\wedge\cdots\wedge dz_m.
\]
Its power transport is
$\eta_0=a(z)^Q\,dw_1\wedge\cdots\wedge dw_m$;
the coefficient $a(z)^Q$ belongs to $R_0$ and is a
unit there.  For $f\in R_1$, the Frobenius trace on
top forms, divided by $\eta_0$, is the $R_0$-linear
functional
\begin{equation}\label{rings:trace-formula}
 \tau(f)=
 \frac{[z_1^{Q-1}\cdots z_m^{Q-1}](f\,a(z))}
      {a(z)^Q}.
\end{equation}
Brackets mean coefficient extraction in the free
$R_0$-basis with all exponents between $0$ and $Q-1$.
This trace is independent of the root coordinates:
it is the finite-duality trace on top forms, or,
equivalently, the functional characterized by the
change-of-variables rule for the residue.  Invariant
volume and the same rule make it translation invariant.

To determine its normalization, make a faithfully
flat extension that supplies a point $b$ of the torsor.
The torsor becomes the finite Honda kernel.  In the
root chart put $\epsilon_i=z_i-b_i$, so that
$\epsilon_i^Q=0$.  Let $u_i$ be normalized group
coordinates at the identity.  Invariance of the
volume gives
\[
 \det\left(\frac{\partial u}{\partial\epsilon}\right)(0)
       =a(b).
\]
In the truncated polynomial ring the socle therefore
transforms as
\[
 u_1^{Q-1}\cdots u_m^{Q-1}
    =a(b)^{Q-1}\epsilon_1^{Q-1}\cdots\epsilon_m^{Q-1}.
\]
Only the linear part of the coordinate change
contributes in the maximal possible total degree;
its action on that socle is the determinant to the
$(Q-1)$-st power.  Formula~\eqref{rings:trace-formula}
consequently sends this group-coordinate socle to
\[
 \frac{a(b)^{Q-1}a(b)}{a(b)^Q}=1.
\]
The normalized invariant Hopf integral has exactly
this value.  The invariant-functional module is a
free line, as is also clear from the perfect
complementary-monomial pairing.  The two functionals
are therefore equal over the whole splitting
algebra.  Faithfully flat descent proves equality
on the original torsor.  In particular no
comparison only on geometric fibers, and no
undetermined invertible function on the base, is
being used.

Identify the root ring with the original ring by
$Q$-powering.  Formula~\eqref{rings:trace-formula}
is then exactly \eqref{rings:stationary-c}.  The
transitivity of finite-duality traces, or direct
iteration of the coefficient-extraction formula,
gives $S^a$ at the $a$-th step.  This proves
stationarity with compatible normalizations.
Under the regular-module identifications in
Lemma~\ref{rings:dual-distribution}, the finite
transfers are the quotient maps
$\Lambda_{q'}\to\Lambda_q$.  Indeed, in normalized
group coordinates the transfer from root exponent
$Q'$ to $Q$ takes the highest monomial to the
highest monomial: extracting the residue classes
of exponents modulo $Q'/Q$ leaves precisely
the exponents $Q-1$.  The invariant-volume
density and its inverse contribute only their
constant terms to these socles.  The transfer
is $T$-equivariant, so its value on this cyclic
generator determines it on the whole regular module.
For $t=1$ the same calculation uses the invariant
forms $dy_i/y_i$ in multiplicative coordinates.
It extracts the zero character on a finite
diagonalizable torsor, so agrees with the
constant-character projection described above.

For the transpose, take first a free compact lattice
over $R_0$ and pair it with the negative principal
parts of its dual tensored with top forms.  The
pairing is
\[
 (f,h\eta)\longmapsto
 \operatorname{res}(fh\eta).
\]
In coordinates, Cartier selects exponents congruent
to $Q-1$ in every variable and divides their
successors by $Q$.  Since $Q$ fixes $k$, its residue
identity is
\begin{equation}\label{rings:transpose-residue}
 \operatorname{res}\bigl(S(f)h\eta\bigr)
       =\operatorname{res}\bigl(fh^Q\eta\bigr).
\end{equation}
This follows as well by applying Cartier to
$fh^Q\eta$ and using its inverse-semilinearity.
Thus the transpose is $h\mapsto h^Q$ in the
stationary volume frame.  Renaming the root
coordinates makes this the root inclusion.  The
form $\eta$ is retained once, as in
\eqref{rings:volume-character}.

For a finite \'etale coefficient algebra the dual
is identified by its finite \'etale trace pairing.
Cartier commutes with \'etale pullback and trace;
this can be checked after a faithfully flat
\'etale splitting, where it is the preceding
formula on each factor.  Hence the transpose is
the actual Frobenius on the finite-cover functions,
transported through their trace dual, and not an
arbitrary semilinear matrix with the same ranks.
The same assertion holds for associated modules
and their morphisms by finite flat descent and
the naturality of the pairing.

Finally, coordinate changes, torsor automorphisms,
and finite-level maps preserve finite duality and
the normalized invariant-volume construction.
Their character changes in source and target agree
because the chosen Frobenius fixes $k$.  This
proves the equivariance assertions, including
those for complete finite diagrams in
Lemma~\ref{rings:coinduction}.  All formulas involve
finite coefficient data and preserve the common-stage
convention of Definition~\ref{rings:bounded-cochains}.
\end{proof}

The last assertion has a useful integral consequence.
On an analytic affinoid chart contained in $x\ne0$,
a finite free coefficient basis and the intrinsic
norm on its finite \'etale cover differ by some
finite factor $C$.  Rooting the actual cover
functions replaces this factor by $C^{1/Q^a}$.
It tends to $1$.  The support comparison will use
this consequence of the explicit transpose, along
with positive lattices constructed in
Lemma~\ref{lem:compact:stable-lattices}; it does
not assign such a norm interpretation to an
arbitrary semilinear endomorphism.
